%!TEX root = main.tex

Our goal is to construct from the type mapping $tm$ and $M$ a mapping $tmap$ that records the typing information for the components in $M$. 
Nevertheless, there is a mismatch between the types in $tm$ and those in $tmap$: While the types in $tm$ are from $\gtypes$, the types in $tmap$ are from $\ftypes$. 
In the sequel, we shall explain how to tackle this mismatch in the construction of $tmap$. 

We shall construct a function $tmap$ from $\range{id}$ to $\ftypes \times \{\duplex, \source, \sink\}$, where in addition to the typing information, the orientations of components are also recorded. The orientations are related to the connect statements. 
There are three types of orientations, \duplex, \source, and \sink. 
In each connect statement $r <= e$, $r$ is a reference and its orientation should be {\duplex} or \sink.
The rational behind the orientations of components and references in $M$ are, 
\begin{itemize}
\item the orientations of wires and registers are $\duplex$, 
\item the orientations of inputs and nodes are \source, 
\item the orientations of outputs are \sink, 
\item the orientation of a reference $t[n]$ are the same as that of $t$, 
\item the orientation of a reference $x.f$ (where $f$ is a field) can be determined from the orientation of $x$ and the flipping status of the field $f$.  
\end{itemize}   


\zhilin{stopped here}


The mapping $tmap$ is constructed by calling the procedure $ports\_stmts\_tmap(pp, ss, tm)$ (see Table~\ref{tab:ports-stmts-tmap}), which then calls the procedures $ports\_tmap(pp, tm)$ and $stmts\_tmap((pmap, pmap), ss, tm)$ respectively. 

Intuitively, $ports\_tmap(pp, tm)$ constructs a mapping from $tm$ and the port declarations $pp$, and $stmts\_tmap((pmap, pmap), ss, tm)$ constructs from $(pmap, pmap)$, the statement sequence $ss$, and $tm$, a pair $(tmap, scope)$. Note that $scope$ is introduced to check that the component references in $ss$ do not violate the scoping constraint introduced by \prettyll{when} statements. For instance, in the following {\hifirrtl} program, there is a scoping error, since $z$ is declared in the scope of the \prettyll{when} statement, but it is referred outside it (Line 7). 
\begin{figure*}[h]
\centering
\begin{tabular}{c}
%[basicstyle=\footnotesize\ttfamily]
\begin{lstlisting}
    input c: UInt<1>
    wire x: SInt<4>
    wire y: SInt<4>
    when c
      wire z: SInt
      z <= x
    z <= y
\end{lstlisting}
\end{tabular}
%\label{fig:example}
\end{figure*}

The procedures $ports\_tmap(pp, tm)$ and $stmts\_tmap((pmap, pmap), ss, tm)$ are defined by the rules in Table~\ref{tab:ports-stmts-tmap}.
To help the readers understand these rules, let us explain the rules ports\_tmap: input and stmts\_tmap: $\neq$nil. 
\begin{itemize}
\item In the rule {ports\_tmap: input}, suppose $pp = input\ x: t\ pp'$. Then $made\_explicit$ is called to instantiate $t$ into $newt$ by $tm$, where all the bit widths are explicit. (The definition of $made\_explicit$ will be given in Table~\ref{tab:made-explicit}.) Let $tmap'$ be the mapping that is obtained by recursively calling $ports\_tmap$ on $(pp', tm)$. Then it is required that $tmap'(id(x)) = none$, that is, $x$ is not declared in $pp'$. Finally, $tmap$ is obtained from $tmap'$ by assigning $(newt, \source)$ to $id(x)$.  
%
\item In the rule {stmts\_tmap: $\neq$nil}, suppose $ss = s\ ss'$. Then $stmt\_tmap(s, tm)$ is called to compute $(tmap'', scope'')$. Moreover, $stmts\_tmap$ is recursively called on $((tmap'', scope''), ss', tm)$ to compute $(tmap', scope')$. 
\end{itemize}
Note that to avoid the tediousness, we omit some rules for $ports\_tmap$ and $stmts\_tmap$ in Table~\ref{tab:ports-stmts-tmap}.  For instance, if $tmap'(id(x)) = none$ in the rule {ports\_tmap: input} is changed to $tmap'(id(x)) \neq none$, then $ports\_tmap(pp, tm) = none$. Intuitively, we omit the rules corresponding to the abnormal situations where the result is $none$. 

\begin{table}[htbp]
  \small
  \begin{gather*}
    \inferii
    {\seqq{ports\_stmts\_tmap(pp, ss, tm) = tmap}}
    {\seqq{ports\_tmap(pp, tm) = pmap}}
    {\seqq{stmts\_tmap((pmap, pmap), ss, tm) = (tmap, \_)}}
   \quad
   \\[1ex]
    \infer[ports\_tmap: input]
    {\seqq{ports\_tmap(pp, tm) = tmap}}
    {
    \begin{array}{c}
    p p= input\ x: t\ pp' \wedge made\_explicit(t, (id(x),0), tm) = (newt, \_) \\
    ports\_tmap(pp', tm) = tmap' \wedge tmap'(id(x)) = none \\
    tmap = add((id(x), (newt, \source)), tmap')
    \end{array}
    }
    \quad
    \\[1ex]
    %
    \infer[ports\_tmap: output]
    {\seqq{ports\_tmap(pp, tm) = tmap}}
    {
    \begin{array}{c}
    p p= output\ x: t\ pp' \wedge made\_explicit(t, (id(x), 0), tm) = (newt, \_) \\
    ports\_tmap(pp', tm) = tmap' \wedge tmap'(id(x)) = none \\
    tmap = add((id(x), (newt, \sink)),tmap')
    \end{array}
    }    
    \quad
   \\[1ex]
    \inferii[stmts\_tmap: nil]
    {\seqq{stmts\_tmap((tmap, scope), ss, tm) = (tmap', scope')}}
    {\seqq{s s = nil}}
    {\seqq{tmap' = tmap \wedge scope' = scope}}
    \quad
   \\[1ex]
    \infer[stmts\_tmap: $\neq$nil]
    {\seqq{stmts\_tmap((tmap, scope), ss, tm) = (tmap', scope')}}
    {
    \begin{array}{c}
      ss = s\ ss' \\
      (tmap'', scope'') = stmt\_tmap(s, tm) \\
      (tmap', scope') = stmts\_tmap((tmap'', scope''), ss', tm)
    \end{array}
    }
  \end{gather*}
%  \vspace{-4mm}
  \caption{$ports\_stmts\_tmap$, $ports\_tmap$, and $stmts\_tmap$}
  \label{tab:ports-stmts-tmap}
\end{table}

Let us also explain some rules of $made\_explicit$. Intuitively, $made\_explicit(ft, id, tm)$ uses $tm$ to determine the implicit widths in $ft$ and obtain $newt$. Moreover, it computes a new identifier $(v, m')$ that is obtained from $id = (v, m)$  by traversing all the ground types in $ft$. 

\zhilin{Explain {made\_explicit: vector}, {type\_eq}, {made\_explicit\_fields: $\neq$nil}: to be done}

\begin{table}[htbp]
  \small
  \begin{gather*}
    \infer[made\_explicit: ground]
    {\seqq{made\_explicit(ft, id, tm) = (newt, newn)}}
    {
    \begin{array}{c}
    ft = (ground, gt) \wedge id = (v, m) \wedge tm(id) = (\_, t)  \wedge type\_eq(gt, t)\\
    newt = t \wedge newn = m + 1
    \end{array}
    }
    \quad
    \\[1ex]
    \infer[made\_explicit: vector]
    {\seqq{made\_explicit(ft, id, tm) = (newt, newn)}}
    {
    \begin{array}{c}
    ft = (vector, t[n]) \wedge id = (v, m) \\
    made\_explicit(t, id, tm) = (t', m') \wedge elsize = m' -  m \\
    \forall 0 \le i < ((n-1)*elsize) \Rightarrow tm(m+i) = tm(m'+i) \\
    newt = (vector, t'[n]) \wedge newn = m + n*elsize 
    \end{array}
    }
   \quad
   \\[1ex]
    \infer[made\_explicit: bundle]
    {\seqq{made\_explicit(ft, id, tm) = (newt, newn)}}
    {
    \begin{array}{c}
    ft = (bundle, ff) \wedge
    made\_explicit\_fields(ff, id, tm) =(newt, newn) 
    \end{array}
    }
    \quad
    \\[1ex]
   \infer[type\_eq]
    {\seqq{type\_eq(cgt, ggt) = res}}
    {
     \begin{array}{c}
     res =
     \big((cgt = (UInt, n) \wedge ggt = (UInt, n)) \vee (cgt = (SInt, n) \wedge ggt = (SInt, n))\ \vee \\
     (cgt = (UIntImp, \_) \wedge ggt = (UIntImp, \_)) \vee (cgt = (SIntImp, \_) \wedge ggt = (SIntImp, \_))\ \vee \\
     (cgt = Clock  \wedge ggt = Clock) \vee (cgt = Reset \wedge ggt = (UInt, 1))\ \vee \\
     (cgt = Reset \wedge ggt = AsyncReset) \vee (cgt = AsyncReset \wedge ggt = AsyncReset) \big)
     \end{array}
    }
    \quad
    \\[1ex]  
    \infer[made\_explicit\_fields: nil]
    {\seqq{made\_explicit\_fields(ff, id, tm) = (newt, newn)}}
    {
    \begin{array}{c}
    ff = nil \wedge id = (v, m) \wedge
    newt = (bundle, nil) \wedge newn = m  
    \end{array}
    }
    \quad
   \\[1ex]
    \infer[made\_explicit\_fields: $\neq$nil]
    {\seqq{made\_explicit\_fields(ff, id, tm) = (newt, newn)}}
    {
    \begin{array}{c}
    ff = (f, b, t, ff') \wedge id = (v, m) \wedge made\_explicit(t, id, tm) = (t', m') \\ 
    made\_explicit\_fields(ff', (v, m'), tm) = (ff'', m'') \\
    newt = (f, b, t', ff'') \wedge newn = m''  
    \end{array}
    }
  \end{gather*}
  \caption{$made\_explicit$, $type\_eq$, and $made\_explicit\_fields$}
  \label{tab:made-explicit}
\end{table}

\begin{table}[htbp]
  \small
  \begin{gather*}
   \infer[size\_of\_ftype: ground]
    {\seqq{size\_of\_ftype(ft) = 1}}
    {\seqq{ft = (ground, gt)}}
    \quad
    \\[1ex]
   \inferii[size\_of\_ftype: vector]
    {\seqq{size\_of\_ftype(ft) = m*n}}
    {\seqq{ft = (vector, ft'[m])}} {\seqq{n = size\_of\_ftype(ft')}}
    \quad
    \\[1ex]
   \inferii[size\_of\_ftype: bundle]
    {\seqq{size\_of\_ftype(ft) = n}}
    {\seqq{ft = (bundle, ff)}} {\seqq{n = size\_of\_fields(ff)}}
    \quad
    \\[1ex]
   \infer[size\_of\_fields: nil]
    {\seqq{size\_of\_fields(ff) = 0}}
    {\seqq{f f = nil}}
    \quad
    \\[1ex]
   \inferiii[size\_of\_fields: $\neq$ nil]
    {\seqq{size\_of\_fields(ff) = m+n}}
    {\seqq{f f = (\_, \_, ft', ff')}} {\seqq{m = size\_of\_ftype(ft')}} {\seqq{n = size\_of\_fields(ff')}}
  \end{gather*}
%  \vspace{-4mm}
  \caption{$size\_of\_ftype$ and $size\_of\_fields$}
  \label{tab:size-list-ftype-fields}
\end{table}

\begin{table}[htbp]
  \small
  \begin{gather*}
    \inferii[stmt\_tmap: skip]
    {\seqq{stmt\_tmap((tmap, scope), s, tm) = (tmap', scope')}}
    {\seqq{s = skip}}
    {\seqq{tmap' = tmap \wedge scope' = scope}}
   \quad
   \\[1ex]
%    \infer[stmt\_tmap: connect-none]
%    {\seqq{stmt\_tmap((tmap, scope), s, vm) = none}}
%    {
%    \begin{array}{c}
%      s = r <= e \\
%      type\_of\_ref(r, scope) = none \vee type\_of\_exp(e, scope) = none
%     \end{array}
%    }
%   \quad
%   \\[1ex]
    \infer[stmt\_tmap: connect]
    {\seqq{stmt\_tmap((tmap, scope), s, tm) = (tmap', scope')}}
    {
    \begin{array}{c}
      s = r <= e \\
      type\_of\_ref(r, scope) \neq none  \\
      type\_of\_exp(e, scope) \neq none \\ 
      (tmap', scope') = (tmap, scope)
     \end{array}
    }
   \quad
   \\[1ex]
%    \infer[stmt\_tmap: invalid-none]
%    {\seqq{stmt\_tmap((tmap, scope), s, vm) = none}}
%    {
%    \begin{array}{c}
%      s = r \mbox{ is invalid} \wedge type\_of\_ref(r, scope) = none
%     \end{array}
%    }
%   \quad
%   \\[1ex]
    \infer[stmt\_tmap: invalid]
    {\seqq{stmt\_tmap((tmap, scope), s, tm) = (tmap', scope')}}
    {
    \begin{array}{c}
      s = r \mbox{ is invalid} \\
      type\_of\_ref(r, scope) \neq none \\
      (tmap', scope') = (tmap, scope)
     \end{array}
    }
   \quad
   \\[1ex]
%    \infer[stmt\_tmap: wire-none]
%    {\seqq{stmt\_tmap((tmap, scope), s, vm) = none}}
%    {
%    \begin{array}{c}
%      s = \mbox{wire } x: t \wedge 
%      (tmap(id(x)) \neq none \vee made\_explicit(t, (id(x),0), vm) = none)
%     \end{array}
%    }
%   \quad
%   \\[1ex]
    \infer[stmt\_tmap: wire]
    {\seqq{stmt\_tmap((tmap, scope), s, tm) = (tmap', scope')}}
    {
    \begin{array}{c}
      s = \mbox{wire } x: t \wedge tmap(id(x)) = none  \\
      made\_explicit(t, (id(x),0), tm) = (newt, \_) \\
      tmap' = add((id(x), (newt, \duplex)), tmap) \\
      scope' = add((id(x), (newt, \duplex)), scope)
     \end{array}
    }
   \quad
   \\[1ex]
%    \infer[stmt\_tmap: node-none]
%    {\seqq{stmt\_tmap((tmap, scope), s, vm) = none}}
%    {
%    \begin{array}{c}
%      s = \mbox{node } x = e \wedge 
%      (tmap(id(x)) \neq none \vee type\_of\_exp(e, scope) = none)
%     \end{array}
%    }
%   \quad
%   \\[1ex]
    \infer[stmt\_tmap: node]
    {\seqq{stmt\_tmap((tmap, scope), s, tm) = (tmap', scope')}}
    {
    \begin{array}{c}
      s = \mbox{node } x = e \wedge tmap(id(x)) = none  \\
      type\_of\_exp(e, scope) = newt \\
      tmap' = add((id(x), (newt, \source)), tmap) \\
      scope' = add((id(x), (newt, \source)), scope)
     \end{array}
    }
   \quad
   \\[1ex]
    \infer[stmt\_tmap: register-asyncreset]
    {\seqq{stmt\_tmap((tmap, scope), s, tm) = (tmap', scope')}}
    {
    \begin{array}{c}
      s = \mbox{reg } x: t \ clk \mbox{ with: reset =>} (rsig, rval)  \\
      tmap(id(x)) = none \wedge type\_of\_exp(clk, scope) \neq none  \\
      made\_explicit(t, (id(x),0), tm) = newt \\
      type\_of\_exp(rsig, scope) = (ground, AsyncReset) \\
      type\_of\_exp(rval, scope) \neq none \\
      tmap' = add((id(x), (newt, \duplex)), tmap) \\
      scope' = add((id(x), (newt, \duplex)), scope)
     \end{array}
    }
   \quad
   \\[1ex]
    \infer[stmt\_tmap: register-reset]
    {\seqq{stmt\_tmap((tmap, scope), s, tm) = (tmap', scope')}}
    {
    \begin{array}{c}
      s = \mbox{reg } x: t \ clk \mbox{ with: reset =>} (rsig, rval)  \\
      tmap(id(x)) = none \wedge type\_of\_exp(clk, scope) \neq none  \\
      made\_explicit(t, (id(x),0), tm) = newt \\
      type\_of\_exp(rsig, scope) = (ground, (UInt, 1)) \\
      scope' = add((id(x), (newt, \duplex)), scope) \\
      type\_of\_exp(rval, scope') \neq none \\
      tmap' = add((id(x), (newt, \duplex)), tmap)
     \end{array}
    }
   \quad
   \\[1ex]
   \infer[stmt\_tmap: when]
    {\seqq{stmt\_tmap((tmap, scope), s, tm) = (tmap', scope')}}
    {
    \begin{array}{c}
      s = \mbox{when } c: sst \mbox{ else }: ssf  \\
      type\_of\_exp(e, scope) \neq none \\
      stmts\_tmap((tmap, scope), sst, tm) = (tmapt, \_) \\
      stmts\_tmap((tmapt, scope), ssf, tm) = (tmapf, \_) \\
      tmap' = tmapf \wedge scope' = scope
     \end{array}
    }
  \end{gather*}
%  \vspace{-4mm}
  \caption{$stmt\_tmap$}
  \label{tab:stmts-tmap}
\end{table}




\hide{
For a statement $\kappa\ v: t$ where $\kappa =\mbox{\prettyll{input},\prettyll{output},\prettyll{wire}}$, 
We define a mapping $tmap$ that assigns to each identifier from $\range{id}$ a type where all widths are explicit.    
For each $v \in \dom(id)$, 
\begin{itemize}
\item if $v$ is declared in a statement $p \equiv \{\mbox{\prettyll{input}} v: t$, then 
\begin{itemize}
\item if $t = \theta$ with $\theta \in \{\mbox{\prettyll{Clock}}, \mbox{\prettyll{Reset}} \mbox{\prettyll{AsyncReset}}\}$ and $tm((id(v),0)) = (\inport, \theta)$, then $tmap(id(v)) = (gt, \theta)$, 
%
\item if $t = \theta$ with $\theta \in \{\mbox{\prettyll{UInt}}, \mbox{\prettyll{SInt}}\}$ and $tm((id(v), 0)) = (\inport, (\theta, n))$, then $tmap(id(v)) = (gt, \theta \langle n \rangle)$, 
%
\item if $t = \theta\langle n \rangle$ with $\theta \in \{\mbox{\prettyll{UInt}}, \mbox{\prettyll{SInt}}\} $ and $tm((id(v),0)) = (\inport, (\theta, n))$, then $tmap(id(v)) = (gt, \theta \langle n \rangle)$, 
%
\item if $t = t'[n]$ and $tm((id(v),n)) = (\inport, )$, then $tmap(id(v)) = (gt, \theta \langle n \rangle)$
\end{itemize}
\item if $v$ is declared in a statement $p \equiv \{\mbox{\prettyll{output}} v: t$, then $tmap(id(v))$ is defined similarly to the case $p \equiv \{\mbox{\prettyll{input}} v: t$, except that $\inport$ is replaced by $\outport$, 
%   
\item if $v$ is declared in a statement $s \equiv \mbox{\prettyll{wire} } v: t$, then 
\begin{itemize}
\item if $v$ is declared in a statement $p \equiv \{\mbox{\prettyll{input}} v: t$, then 
\begin{itemize}
\item if $t = \theta$ with $\theta \in \{\mbox{\prettyll{Clock}}, \mbox{\prettyll{Reset}} \mbox{\prettyll{AsyncReset}}\}$ and $tm(id(v)) = (\inport, \theta)$, then $tmap(id(v)) = (gt, \theta)$, 
%
\item if $t = \theta$ with $\theta \in \{\mbox{\prettyll{UInt}}, \mbox{\prettyll{SInt}}\}$ and $tm(id(v)) = (\inport, (\theta, n))$, then $tmap(id(v)) = (gt, \theta \langle n \rangle)$. 
\end{itemize}
\end{itemize}
\end{itemize}
}




%A HiFIRRTL program $P$ is mapped to \emph{a set of} module graphs. %, where the box 
%We define $\mathcal{G}(P)$.  % G(P): the set of module graphs of  P


%We assume a component environment $ce$ which extracts the type information from $P$ %the sequence $ss$ of statements 
%and 
%the connect mapping $cm$ extracts for each component the expression that is connected to the component according to the last-connect semantics. 

%Recall that a circuit defined by $P$ consists of multiple modules. 
%
%\begin{itemize}
%\end{itemize}

%	\item tg
%\end{itemize} 


%A HiFIRRTL module $M$ will be first transformed into a circuit diagram $D_M = (ss, ce, cm)$, where $ss$ is the sequence of statements in $M$, $ce$ is the component environment obtained from the variable declarations in $M$, possibly with the types of some variables missing or partial, 

%and $cm$ is the trivial mapping that maps each component to $\bot$ (Intuitively, this means that the connects of the components have not be determined yet). The missing or partial  types in $ce$ will be completed by the compilation passes during the transformation of $M$ into a LoFIRRTL module. Moreover, the connect mapping $cm$ will be updated by the ExpandWhens compilation pass. Furthermore, the statement sequence $ss$ will be updated by ExpandConnects, ExpandWhens, LowerTypes, and RemoveReset, while they keep unchanged during the other compilation passes.  



%\subsection{Behavioral semantics of modular graphs}
